Training data deduplication is widely believed to improve language model quality,
yet when we compare Pythia models trained on the Pile versus its deduplicated variant
at token-count-matched checkpoints, deduplication makes MMLU accuracy \textit{significantly
worse} while improving HellaSwag and ARC-Challenge. We show this divergence is not
paradoxical --- it is a contamination-correction signature: the direction and magnitude
of each benchmark's accuracy change under deduplication correlates positively with the
benchmark's estimated n-gram overlap with the Pile training corpus
(Pearson $r = 0.632$, $p = 0.0086$ across 16 observations spanning 4 benchmarks and
4 model sizes 160M--6.9B). Benchmarks most contaminated in the original corpus show
the largest accuracy reductions under deduplication, while low-contamination benchmarks
are unaffected or improved. We additionally demonstrate that token-count matching ---
not step-matching --- is the correct confound-control methodology for such comparisons,
recovering a $\Delta r = 0.093$ stronger contamination signal than step-matched analyses.
Our findings reframe deduplication's benchmark effects from a uniform quality improvement
to a selective contamination correction, with direct implications for how practitioners
should interpret benchmark regressions after corpus deduplication.
